\section{Rules and stateless learners}
\label{ch:stateless}

The families without state come first: a causal rule on the market's own price, or a learner that
reads a few dozen numbers computed from the clean data at the sample time and nothing else.
It is the family that tells where the information is.

\subsection{Causal rules}
\label{sec:stateless:rules}


\begin{table}[t]
\centering
\caption{Causal rules on the validation day, balanced accuracy, all samples and
deduplicated samples (one per market per ten minutes).  Momentum-$H$ and Mean
Reversion-$H$ use the change over the past $H$; Momentum last-tick uses the sign of the
last mid change.  At 1 min and 15 min only Mean Reversion was evaluated, on the samples of
the short-horizon runs.}
\label{tab:exp:baselines}
\small
\begin{fittable}\begin{tabular}{lcccccccccc}
\toprule
& \multicolumn{2}{c}{1 min} & \multicolumn{2}{c}{15 min} & \multicolumn{2}{c}{1 h} & \multicolumn{2}{c}{6 h} & \multicolumn{2}{c}{12 h} \\
\cmidrule(lr){2-3}\cmidrule(lr){4-5}\cmidrule(lr){6-7}\cmidrule(lr){8-9}\cmidrule(lr){10-11}
Rule & all & dedup & all & dedup & all & dedup & all & dedup & all & dedup \\
\midrule
Constant, prior, uniform  & 0.333 & 0.333 & 0.333 & 0.333 & 0.333 & 0.333 & 0.333 & 0.333 & 0.333 & 0.333 \\
Momentum last-tick        & --- & --- & --- & --- & 0.375 & 0.387 & 0.365 & 0.374 & 0.362 & 0.370 \\
Momentum-$H$              & --- & --- & --- & --- & 0.472 & 0.445 & 0.436 & 0.422 & 0.424 & 0.420 \\
Mean Reversion-$H$        & \textbf{0.536} & \textbf{0.501} & \textbf{0.537} & \textbf{0.525} & \textbf{0.527} & \textbf{0.506} & \textbf{0.501} & \textbf{0.488} & \textbf{0.507} & \textbf{0.500} \\
\bottomrule
\end{tabular}\end{fittable}
\end{table}

Mean Reversion is above $0.5$ at every horizon and it is the rule to beat
(\cref{tab:exp:baselines}).  Momentum is far below it and, on the sign of the last tick,
close to the constant rule: on this data the sign of the last move is \emph{anti}-informative,
and the whole of the simple predictability is a reversal.  The Mean Reversion values quoted
in later sections differ in the third decimal from these because every comparison is made on
the samples of the run it belongs to (\cref{sec:models:protocol}).

\subsubsection{Why mean reversion is strong}
\label{sec:analysis:extreme}


A Polymarket mid price is a bounded, discrete process observed through a thin book.
Most of its changes are one-tick moves produced by an order arriving at or leaving the
best quote; they are not information and they revert.  A market that has not moved in
the last hour is usually illiquid or settled near a boundary and will not move in the
next hour either.  Mean Reversion-$H$ encodes exactly these two regularities, and it is why a
rule that uses one number is above $0.5$ balanced accuracy.  Its strength is largest at
extreme prices (below $0.05$ or above $0.95$), where the process is almost purely a
one-tick oscillation, and it is there that every learned model is weakest relative to it.

\subsection{Trees on hand-made features: the no-state control}
\label{sec:stateless:control}
\label{sec:exp:nostate}

The sixty-one features of \cref{sec:prep:features} were given to LightGBM one group at a
time, on the validation samples of the best one-hour stateful model (the time-first model
of \cref{sec:stateful:encoders}), so that the two can be compared on the same samples.
\begin{table}[t]
\centering
\caption{Gradient boosting without state against the deep models, 1h, on the same
validation samples (balanced accuracy; last column: difference from the time-first model
in points with 95\% hour-block bootstrap intervals).}
\label{tab:exp:nostate}
\small
\begin{fittable}\begin{tabular}{lrccccc}
\toprule
Model & \# & All & Dedup. & Extreme & Central & $\Delta$ vs time-first \\
\midrule
Mean Reversion                 &    & 0.5311 & 0.5180 & 0.5230 & 0.5348 & $-3.32$ \\
GRU $K{=}10$ model             &    & 0.5547 & 0.5382 & 0.5385 & 0.5602 & $-0.97$ [$-1.28$, $-0.63$] \\
Time-first $K{=}64$ model      &    & 0.5643 & 0.5410 & 0.5377 & 0.5748 & \\
\midrule
G1 price history               & 12 & 0.5617 & 0.5509 & 0.5449 & 0.5685 & $-0.26$ [$-0.67$, $+0.17$] \\
G2 + book shape                & 26 & 0.5831 & 0.5718 & 0.5620 & 0.5918 & $+1.87$ [$+1.55$, $+2.19$] \\
G3 + activity and clock        & 47 & 0.5929 & 0.5757 & 0.5722 & 0.6015 & $+2.86$ [$+2.54$, $+3.16$] \\
G4 + on-chain trades           & 56 & 0.5997 & 0.5812 & 0.5788 & 0.6083 & $+3.54$ [$+3.22$, $+3.88$] \\
G5 + static fields             & 61 & \textbf{0.6015} & \textbf{0.5827} & \textbf{0.5804} & \textbf{0.6107} & $+3.72$ [$+3.33$, $+4.12$] \\
\bottomrule
\end{tabular}\end{fittable}
\end{table}

Price history alone reproduces the deep model.  The book shape adds two points, the
activity clock one more, the trades two thirds of a point and the static fields a fifth;
the full model is $3.7$ points above the best deep model on the full stream, $4.2$ on the
deduplicated samples, and ahead in every one of the 24 validation hours and in every
market category.  The gain importance is dominated by the time since the last mid change
(32\%), the number of mid changes in the last hour (14\%) and the realised volatility
(10\%); the level-1 imbalance and the category get 1\% each and every order-flow
accumulator of the tensor less than that.  Univariately, on the samples that moved, the
sign of ``mid minus its one-hour mean'' gives the direction 57.3\% of the time, the
level-1 imbalance 55.7\%, the microprice 55.5\%, the on-chain flow 52.6\% (on the fifth
of the samples with trades), the 10-level imbalance 50.8\% and the tensor's all-level
order-flow imbalance 50.0--50.5\%; for separating \textsc{const} from a move, the time
since the last mid change has an AUC of 0.85.

\subsubsection{The other horizons}
\label{sec:stateless:horizons}

\begin{figure}[t]
\centering
\begin{tikzpicture}
\begin{axis}[thesis, ybar, bar width=7pt, ymin=0.45, ymax=0.65, ylabel={balanced accuracy, all samples},
  symbolic x coords={1 min, 15 min, 1 h, 6 h, 12 h}, xtick=data, enlarge x limits=0.14,
  legend columns=4, legend style={at={(0.5,1.03)}, anchor=south}, ytick={0.45,0.50,0.55,0.60,0.65}]
\addplot[fill=vizD, draw=white, line width=0.5pt] coordinates {(1 min,0.5363) (15 min,0.5373) (1 h,0.5311) (6 h,0.5069) (12 h,0.5071)};
\addplot[fill=vizA, draw=white, line width=0.5pt] coordinates {(1 min,0.5022) (15 min,0.5689) (1 h,0.5617) (6 h,0.5380) (12 h,0.5446)};
\addplot[fill=vizB, draw=white, line width=0.5pt] coordinates {(1 min,0.6065) (15 min,0.6116) (1 h,0.5929) (6 h,0.5606) (12 h,0.5616)};
\addplot[fill=vizC, draw=white, line width=0.5pt] coordinates {(1 min,0.6086) (15 min,0.6160) (1 h,0.6015) (6 h,0.5716) (12 h,0.5715)};
\legend{Mean Reversion, G1 price history, G3 + book + activity, G5 all 61}
\end{axis}
\end{tikzpicture}

\caption{The stateless control by feature group at every horizon, against Mean Reversion
(all samples).  Price history alone (G1) is below the rule at one minute and above the
stateful model at six and twelve hours; the activity clock (G3) carries most of the gain at
the short horizons.}
\label{fig:stateless:groups}
\end{figure}

The same control was run at the other horizons (\cref{tab:exp:nostate:horizons,fig:stateless:groups}): at
6h and 12h on the exact validation samples of the deep models, at 1 minute, where no deep
model has saved predictions, on one tensor row in twenty of the validation day labelled
like the training rows, against Mean Reversion alone.  At 6h and 12h the full model is
$6.4$--$6.5$ points above the rule: $3.2$ above the time-first model at 6h and $8.0$ above
the reference at 12h, where the deep model is below the rule.  Unlike at 1h, price history
alone is already clearly above the deep model at these horizons ($+1.7$ at 6h, $+5.3$ at
12h): the five returns read by trees together with the price level carry most of the
long-horizon signal, and the 12h return is the strongest single feature (reversal sign,
$57.9\%$ directional hit).  The activity clock keeps a third of the gain importance, and
the time to the end of the market, its age and its duration rise to about $9\%$ together.
At 1 minute the direction of a move is nearly unpredictable (the best univariate hit
rate is $54.1\%$, the level-1 imbalance), and the whole gain is in \emph{whether} the mid
will move: the number of mid changes in the last hour separates \textsc{const} from a
move with an AUC of $0.94$, the time since the last change $0.91$, the event count of the
last minute $0.89$, and the activity group alone is $7.2$ points above the rule.  Price
history alone is in fact \emph{below} the rule at 1 minute on the deduplicated ($-6.1$)
and extreme ($-7.1$) samples: trees on returns learn the move-or-not prior of the full
stream, which is wrong on independent samples.  The ordering of the groups is the same
at every horizon, activity clock first, then book shape, then price history, with trades
and static fields marginal; the tensor's order-flow accumulators stay below $1\%$ of the
importance everywhere.

\begin{table}[t]
\centering
\caption{The no-state control at the other horizons, all samples and deduplicated samples
(balanced accuracy; $\Delta$ in points with 95\% hour-block bootstrap intervals).  At 1
minute the samples are one tensor row in twenty of the validation day and the only
comparator is Mean Reversion.}
\label{tab:exp:nostate:horizons}
\small
\begin{fittable}\begin{tabular}{llcccccc}
\toprule
Horizon & Subset & MR & Best deep & G1 & G3 & G5 & $\Delta$ G5 $-$ MR [95\% CI] \\
\midrule
6h    & All   & 0.5069 & 0.5396 (TF) & 0.5380 & 0.5606 & \textbf{0.5716} & $+6.47$ [$+6.02$, $+6.91$] \\
      & Dedup & 0.5002 & 0.5370 (TF) & 0.5360 & 0.5588 & \textbf{0.5665} & $+6.63$ \\
12h   & All   & 0.5071 & 0.4914 (GRU) & 0.5446 & 0.5616 & \textbf{0.5715} & $+6.44$ [$+5.53$, $+7.33$] \\
      & Dedup & 0.4995 & 0.5168 (GRU) & 0.5427 & 0.5598 & \textbf{0.5680} & $+6.85$ \\
1 min & All   & 0.5862 & --- & 0.5901 & 0.6579 & \textbf{0.6594} & $+7.32$ [$+6.22$, $+8.65$] \\
      & Dedup & 0.5244 & --- & 0.4636 & 0.5604 & \textbf{0.5610} & $+3.67$ \\
\bottomrule
\end{tabular}\end{fittable}
\end{table}
Two caveats.  The boosted model optimises balanced accuracy directly through its class
weights, whereas the deep model is trained with the focal loss and predicts \textsc{up}
for 15\% of the samples against 21.5\% in the labels, so part of the gap may be
calibration (the deep model's probabilities were not saved, so a prior correction could not
be tried).  And the control was trained on a tenth of five days; the deep model on seven
full days.  Neither caveat is large enough to change the sign.

\subsubsection{Where the information is}
\label{sec:analysis:nostate}

The control of \cref{sec:exp:nostate} is the most informative negative result of the
thesis.  A model with no memory, reading a few dozen numbers that anyone would compute
from the clean data, beats the best recurrent model by almost four points.  Three things
follow.  First, the recurrent state is not where the information lives: the price history
alone (five returns, the mid and the hour, fed to trees instead of to a linear direct path)
already matches the deep model, so what the state adds on top of the heuristic is small.
Second, the information the deep model misses is not hidden in the order book; it is the
\emph{activity clock}.  The single best predictor of whether a price will move in the
next hour is how long ago it last moved, followed by how many times it moved in the last
hour; neither is a column of the tensor, and the event-clock window (\cref{sec:data:tensors})
hides them, because a window of twenty events can be ten seconds or ten minutes old and
the model is never told which.  Third, the way the book is compressed matters more than
the architecture that reads it, which is the lesson of \textcite{kolm2023dofi} and
\textcite{lucchese2024predictability} reproduced here: the direction is in the first
level (imbalance, microprice) and in the position of the mid relative to its recent mean,
not in ten levels of depth nor in an order-flow accumulator that sums changes at levels the
mid will never reach.

The other horizons sharpen the picture rather than change it.  At 6h and 12h the
control is six and a half points above the rule and the price history alone already
beats the deep model, so the long-horizon signal is a slow reversal of the 6h and 12h
returns, modulated by the price level and by where the market is in its life (time to
the end, age), which trees read directly and the direct path, linear in the returns,
cannot bend to the level.  At 1 minute nothing predicts the direction, every feature is
within four points of a coin toss on the samples that moved, and the whole task is to
predict \emph{activity}: the gain of the control is the activity clock alone, and
returns without it are worse than the rule on independent samples.  The horizon thus
changes what the task is, from ``will the price move in the next minute'' to ``which
way will it revert over the next half day'', and the same architecture cannot be the
right one at both ends.

\subsection{Neural tabular learners: trees against networks}
\label{sec:stateless:tabular}

Whether a gap between the trees and a deep model is the learner or the training regime can be
settled on the control's table alone.  A wider tabular network (three layers of 512 units on the piecewise-linear encoding,
shuffled mini-batches, a one-cycle schedule over sixteen epochs with early stopping) reaches
0.5983 with a single seed and 0.5991 with three; the parameter-efficient ensemble TabM
\cite{gorishniy2025tabm}, sixteen members sharing one backbone, reaches 0.6012, 0.6015 and 0.5996
with three seeds, indistinguishable from LightGBM ($-0.04$ [$-0.21$, +0.11] for the seed average),
and its combination with the trees adds nothing (+0.07): the two learners know the same things.
The claim that trees beat networks on this table does not survive a properly trained network; the
claim that the online, time-ordered regime of the stateful model costs between one and a half and
two points does (one pass 0.5842, two passes 0.5907, tabular 0.6011; \cref{sec:stateful:regime}).

At the short horizons the same learner on the same table gives the same answer
(\cref{ch:horizons}): TabM is $+0.33$ over the calibrated trees at one minute and $+0.22$ at
fifteen, and their combination adds nothing beyond TabM at either.  Trees and a properly
trained network know the same things on this table.

\subsection{What a stateless model cannot see}
\label{sec:stateless:limits}

A stateless learner reads a summary of the past at the sample time: how long since the last
move, how many moves in the last hour, the shape of the book now, the returns over fixed
lookbacks.  What it cannot read is the \emph{sequence} of book events, the order in which
depth appeared and disappeared, the pattern of the last sixty seconds of a quiet market.
Whether that sequence carries information beyond its summary is exactly what the stateful
family measures, and the answer of \cref{ch:stateful,ch:horizons} is: not at one hour,
where the summary is complete, and a little at one minute, where it is not.

